
\section{\trChernClass}
\label{sec:appB-chern-classes}

\emph{第二\trCohomology{}群。}

我们曾经提到过复\trVectorBundle{}的第一\trChernClass{}，尽管我们对它用得不多。它居于\trCohomology{}群 $H^2(W;\ZZ)$ 中——例如，可以像本书第一部分中那样用\trMorseComplex{}得到。

\emph{第一\trChernClass。}

定义它最经济的方式，是取这个\trFiberBundle{}的最大外幂，从而得到一个复线丛；而第一\trChernClass{}就是后者的\trEuler{}类：
\[
c_1(E) = c_1(\Lambda^n E) = e(\Lambda^n(E)) \in H^2(W;\ZZ)
\]
其中 $E$ 是空间 $W$ 上\trRank{}为 $n$ 的\trFiberBundle。这些概念参见 [56]。

我们曾经用到过一个\trSymplectic\trManifold{}的\trTangentBundle{} $TW$ 的第一\trChernClass{} $c_1(TW)$——借由一个被\trSymplecticForm{}所\trCalibrated{}的\trAlmostComplex{}结构，把它看作一个复\trVectorBundle。让我们注意到，这个类是\trWellDefined{}的，理由很简单：这类\trAlmostComplex{}结构所成的集合是\trContractible{}的（从而特别地是\trConnected{}的），而\trChernClass{} $c_1$ 居于离散空间 $H^2(W;\ZZ)$ 之中。
